\section{Consignes de codage}

Ces règles s'appliquent à \textbf{tous} les projets de ce dossier. Elles
seront vérifiées lors de la relecture du code et discutées en
soutenance.

\subsection*{Pas de variable globale non justifiée}
Toute donnée nécessaire à une fonction doit lui être transmise par
paramètre (ou renvoyée par valeur de retour), pas via une variable
globale. Une variable globale ne se justifie que dans des cas rares et
documentés (ex.\ une constante de configuration) -- et dans ce cas, le
commentaire doit expliquer \textbf{pourquoi} elle ne peut pas être un
paramètre.

\subsection*{Noms de variables, fonctions et tableaux (inspiré de la norme 42)}
Un nom doit permettre de comprendre son rôle sans avoir à lire le reste
du code. Pas de préfixe de type façon notation hongroise (\texttt{l\_},
\texttt{tab\_}, \texttt{g\_}...) : le rôle se lit dans le nom, pas dans
un préfixe.

\begin{itemize}
  \item \textbf{snake\_case partout}, jamais de camelCase : \texttt{nb\_joueurs}, pas \texttt{nbJoueurs}.
  \item \textbf{Variables} : nom explicite et complet. Préférer \texttt{nb\_joueurs} à \texttt{n}, \texttt{indice\_carte\_tiree} à \texttt{i2}. Seuls \texttt{i}, \texttt{j}, \texttt{k} sont tolérés, uniquement comme compteur de boucle local et de portée très courte (ex.\ \texttt{for (int i = 0; ...)}).
  \item \textbf{Booléens} : préfixés par \texttt{est\_}, \texttt{a\_} ou \texttt{peut\_} pour que leur lecture soit une question/affirmation : \texttt{est\_gagnant}, \texttt{a\_rejoue}, \texttt{peut\_tirer}.
  \item \textbf{Tableaux} : un nom explicite, au pluriel si possible, décrivant ce qu'il contient -- pas de préfixe \texttt{tab\_} : \texttt{scores}, \texttt{cartes\_tirees}, \texttt{grille}.
  \item \textbf{Fonctions} : \texttt{verbe\_nom}, une action claire par fonction : \texttt{calculer\_score()}, \texttt{afficher\_plateau()}. Une fonction qui renvoie un booléen se lit comme une question : \texttt{est\_victoire()}, \texttt{peut\_jouer()}.
  \item \textbf{Constantes et macros} : \texttt{UPPER\_SNAKE\_CASE} : \texttt{NB\_JOUEURS\_MAX}, \texttt{TAILLE\_GRILLE}.
\end{itemize}

\subsection*{Indentations valides}
Chaque niveau d'imbrication (fonction, boucle, condition) doit être
indenté de façon cohérente et visible, de sorte que la structure du
code soit lisible sans avoir à compter les accolades. Un éditeur de
code peut généralement réindenter automatiquement un fichier -- prendre
l'habitude de le faire avant de rendre son travail.
